% ---------------------------------------------------------------------------
% The variance-gradient proposition: unifies the criterion's three gates as the
% three terms of the within-group reward variance. Drop into the criterion section.
% ---------------------------------------------------------------------------
\subsection{Why the criterion holds: group-relative gradient is potential variance}
\label{sec:theorem}

The three gates are not independent heuristics; they are the three terms of a single
variance. We make this exact for group-relative policy optimization.

\paragraph{Setup.} For a prompt, draw a group of $G$ rollouts $\tau_1,\dots,\tau_G$ from
$\pi_\theta$. Let each rollout's scalar reward decompose as
\[
  R_i \;=\; O_i \;+\; \beta\,\Phi_i,
\]
where $O_i\in\{0,1\}$ is the terminal outcome, $\Phi_i\in\mathbb{R}$ is a verifiable
\emph{potential} summarising rollout $i$'s process (goals discharged, stages completed,
fraction of tests passing), and $\beta$ weights the process channel. Group-relative
optimisation (Dr.\,GRPO mean-centring) assigns advantage $A_i = R_i - \bar R$ with
$\bar R=\frac1G\sum_j R_j$, and forms the estimator
$g=\frac1G\sum_i A_i\,\nabla_\theta\log\pi_\theta(\tau_i)$.

\paragraph{The advantage energy is the within-group reward variance.}
The magnitude of the group's update is carried by the squared advantage weights, and
\[
  \frac1G\sum_{i=1}^{G} A_i^2 \;=\; \frac1G\sum_{i=1}^{G}\big(R_i-\bar R\big)^2
  \;=\; \operatorname{Var}_G(R).
\]
Decomposing $R=O+\beta\Phi$,
\begin{equation}
  \boxed{\;\operatorname{Var}_G(R) \;=\; \underbrace{\operatorname{Var}_G(O)}_{\text{outcome}}
  \;+\; \beta^2\underbrace{\operatorname{Var}_G(\Phi)}_{\text{potential}}
  \;+\; 2\beta\underbrace{\operatorname{Cov}_G(O,\Phi)}_{\text{alignment}}\;}
  \label{eq:vardecomp}
\end{equation}

\begin{proposition}[Where the gradient comes from]
\label{prop:variance}
A group contributes zero to the policy gradient iff $\operatorname{Var}_G(R)=0$, and the
magnitude of its contribution is governed by \eqref{eq:vardecomp}. Consequently the
criterion's three gates are the three terms of \eqref{eq:vardecomp}:
\begin{enumerate}[leftmargin=1.6em,itemsep=1pt,topsep=2pt,label=\textnormal{(\roman*)}]
\item \textbf{All-fail blindness (outcome term).} On an all-fail group $O_i\equiv 0$, so
$\operatorname{Var}_G(O)=\operatorname{Cov}_G(O,\Phi)=0$ and the energy reduces to
$\beta^2\operatorname{Var}_G(\Phi)$. Outcome-only RL ($\beta=0$) yields exactly zero: the
dead update. The \emph{entire} learnable signal on the hardest groups is the potential's
within-group variance.
\item \textbf{Finer \& reachable $\Leftrightarrow$ gradient (potential term).} On all-fail
groups the gradient is non-zero iff $\operatorname{Var}_G(\Phi)>0$, i.e.\ $\Phi$ takes at
least two distinct values among the \emph{sampled} rollouts. ``Finer than the outcome''
is what makes more than one value \emph{possible}; ``reachable'' is what makes more than
one value \emph{actual}. A potential equal to the binary outcome, or one whose intermediate
values are never reached ($\Phi_i\equiv 0$), has zero variance and is centred out---no
gradient, gameable or not.
\item \textbf{Un-gameable $\Leftrightarrow$ the gradient points at success (alignment
term).} The process channel moves $\theta$ along the $\Phi$-weighted score; whether this
helps the \emph{outcome} is governed by the sign of $\operatorname{Cov}_G(O,\Phi)$ under
training. Un-gameability---the cheapest $\Phi$-maximising policy also maximises $O$---is
exactly the condition that keeps this covariance non-negative as the policy improves. A
gameable $\Phi$ admits a policy with $\mathbb{E}[\Phi]\!\uparrow$ while
$\mathbb{E}[O]\!\to\!0$ ($\operatorname{Cov}<0$): the compliance attractor, in which the
potential variance \eqref{eq:vardecomp} is large but points away from success.
\end{enumerate}
\end{proposition}

Parts (i)--(ii) are exact algebra; part (iii) is the dynamical reading---the only place the
\emph{training trajectory} (how $\operatorname{Cov}_G$ evolves) enters, and where the
empirical un-gameability sweep (\S\ref{sec:law}) supplies what the static identity cannot.
The proposition turns the criterion into a statement about one quantity, the within-group
variance of a verifiable potential: \emph{make it positive} (finer, reachable) \emph{and
keep it aligned} (un-gameable), and a dense process reward is exactly the gradient that
outcome-only RL lacks on the groups that matter.
